% Part: first-order-logic
% Chapter: proof-systems
% Section: introduction

\documentclass[../../../include/open-logic-section]{subfiles}

\begin{document}

\iftag{FOL}
      {\olfileid{fol}{prf}{int}}
      {\olfileid{pl}{prf}{int}}

\olsection{مقدمة}

للمنطق في العادة دلالة ونظام !!a{derivation}. أما الدلالة ففيها بحث الصدق وقابلية الإرضاء والصلاحية والاستلزام. وأما أنظمة ال!!{derivation} فغايتها أن يُثبت الاستلزام والصلاحية بطريق تركيبي محض. ومعنى ذلك أن ال!!{derivation} كائن تركيبي منتهٍ: يكون في الغالب متتالية، أو ترتيبًا منتهيًا آخر، من ال!!{sentence}s أو ال!!{formula}s. ومزية أنظمة ال!!{derivation} المحكمة أن صحة ما يُعرض من متتاليات ال!!{sentence}s أو ال!!{formula}s وترتيباتها قابلة للفحص الآلي.

وأبسط أنظمة ال!!{derivation} لمنطق الرتبة الأولى، وأسبقها في التاريخ، الأنظمة \emph{البديهية}. ففيها تُعد متتالية ال!!{formula}s
!!a{derivation} إذا كانت كل !!{formula} فيها إما بديهية من مجموعة ثابتة، وإما نتيجة لِما تقدمها من !!{formula}s بقاعدة من عدد ثابت من قواعد الاستدلال. وكون
!!a{formula} بديهية، أو كونها مستنتجة على وجه صحيح من !!{formula}s أخرى بقاعدة منها، أمر يُتحقق منه آليًا. ويسهل وصف أنظمة ال!!{derivation} البديهية وبحثها فيما وراء النظرية؛ لكن ال!!{derivation}s التي تنشأ فيها صعبة الإنشاء والقراءة والفهم.

ولأجل تيسير إنشاء ال!!{derivation}s أو فهم ال!!{derivation}s التامة وُضعت أنظمة
!!{derivation} أخرى: منها الاستنتاج الطبيعي، وأشجار الصدق، وهي المسماة أيضًا براهين الجداول الدلالية، وحساب التتابعيات. ومنها أنظمة
!!{derivation} روعي في وضعها العمل الآلي خاصة؛ فطريقة الحَلّ سهلة التنفيذ في البرمجيات، وإن كانت !!{derivation}sها تكاد تستعصي على الفهم. ومعظم أنظمة ال!!{derivation} هذه تجعل ال!!{derivation}s أشجارًا من ال!!{formula}s، لا متتاليات. وبذلك يتبيّن في
!!a{derivation} أيُّ الأجزاء مبني على أيِّها.

تختلف إذن أنظمة ال!!{derivation} للمنطق الواحد، كمنطق الرتبة الأولى، في بيان معنى أن تكون
!!a{sentence} \emph{مبرهنة}، ومعنى أن تكون
!!a{sentence} !!{derivable} من غيرها. وسواء استُعملت
!!{derivation}s بديهية، أم استنتاجات طبيعية، أم
!!{derivation}s تتابعية، أم أشجار صدق، أم تفنيدات بالحَلّ، فالمقصود أن تطابق هذه العلاقات الصلاحية والاستلزام الدلاليين. نكتب
$\Proves !A$ لقولنا «$!A$ مبرهنة»، ونكتب
«$\OLId{greek-variable}{Gamma} \Proves !A$» لقولنا «$!A$ !!{derivable} من~$\OLId{greek-variable}{Gamma}$».
فمهما يكن تعريف~$\Proves$، نطلب مطابقته لـ~$\Entails$ على الوجه الآتي:
\begin{enumerate}
\item $\Proves !A$ إذا وفقط إذا كان $\Entails !A$
\item $\OLId{greek-variable}{Gamma} \Proves !A$ إذا وفقط إذا كان $\OLId{greek-variable}{Gamma} \Entails !A$
\end{enumerate}
اتجاه «فقط إذا» هو \emph{السلامة}: يكون نظام
!!^a{derivation} سليمًا إذا كانت !!{derivability} تستلزم الاستلزام الدلالي أو الصلاحية. ولا بد أن يكون كل نظام
!!{derivation} معتبر سليمًا؛ إذ لا تنفع أنظمة ال!!{derivation} غير السليمة أصلًا. فإن الغاية من
!!a{derivation} أن يكون ضمانًا تركيبيًا للصلاحية أو الاستلزام. وسنقيم برهان السلامة لكل ما نعرضه من أنظمة ال!!{derivation}.

وأما الاتجاه المقابل، أي «إذا»، فاسمه \emph{الاكتمال}. نظام ال!!{derivation} الكامل يفي بإثبات أن $!A$ مبرهنة متى كانت $!A$ صالحة، وبإثبات
$\OLId{greek-variable}{Gamma} \Proves !A$ متى كان $\OLId{greek-variable}{Gamma} \Entails !A$.
والاكتمال أصعب إثباتًا، بل لبعض الأنساق المنطقية لا يوجد نظام
!!{derivation} كامل. وليس منطق الرتبة الأولى كذلك؛ فقد أثبت كورت غودل أول مرة اكتمال نظام
!!a{derivation} له في أطروحته سنة \OLClassicalDigits{1929}.

ويتصل بأنظمة ال!!{derivation} معنى \emph{الاتساق}. فمجموعة ال!!{sentence}s غير متسقة إذا أمكن
!!{derive} أي شيء منها، وإلا فهي متسقة. وعدم الاتساق نظير عدم قابلية الإرضاء في الجانب التركيبي: فالمجموعات غير المتسقة من ال!!{sentence}s لا تصلح نظريات جيدة، كما لا تصلح المجموعات غير القابلة للإرضاء؛ إذ الخلل في كلتيهما أساسي. وليس يلزم أن تكون مجموعة ال!!{sentence}s المتسقة صادقة أو نافعة؛ لكنها تجتاز هذا الحد الأدنى من الصلاح للاستعمال المنطقي.
وقد يتغير تعريف الاتساق الدقيق بتغير نظام ال!!{derivation} المستعمل لمجموعات ال!!{sentence}s. غير أن المطلوب فيه، كما في~$\Proves$، موافقة النظير الدلالي، أعني قابلية الإرضاء: تكون~$\OLId{greek-variable}{Gamma}$ متسقة إذا وفقط إذا أمكن إرضاؤها. وفي هذه العبارة يدل اتجاه «فقط إذا» على الاكتمال، إذ يلزم من الاتساق إمكان الإرضاء؛ ويدل اتجاه «إذا» على السلامة، إذ يلزم من إمكان الإرضاء الاتساق. وهاتان الصيغتان للسلامة والاكتمال مكافئتان للصيغتين السابقتين في منطق الرتبة الأولى الكلاسيكي.

\end{document}
